\section{Agency, Ownership, and Self-Location}
\label{sec:background_agency_ownership_self_location}

The sense of embodiment is commonly discussed through three partially separable components: sense of agency, body ownership, and self-location. These components are useful for this thesis because the two rowing representations may not affect all aspects of embodiment in the same way. A representation that supports control may not necessarily feel most like rowing, and a representation that looks most like rowing may not necessarily provide the strongest correspondence with the user's felt body.

Sense of agency refers to the experience of causing or controlling an action. In virtual environments, agency is often supported when a virtual body or object moves in a way that matches the user's own intentions and movements. When a user moves their hand and immediately sees a virtual hand perform the same movement, the visual feedback can support the impression that the virtual action is their own. Conversely, delays, spatial offsets, or unexpected transformations can weaken the sense that the observed action is under direct control.
% TODO cite: sense of agency in VR / motor control / intentional binding if relevant.
% TODO cite: agency and visuomotor congruence source.
% TODO cite: latency thresholds affecting agency if used, e.g. studies around delays >150 ms.

For VR rowing, agency is especially important because the activity is active, continuous, and rhythmic. The user repeatedly initiates and controls a stroke cycle. If the virtual hands, handle, or rowing body visibly follow the user's real movement, agency may be reinforced throughout the stroke. If the virtual action transforms the user's movement into a different oar or boat motion, agency may still be present, but it may depend on whether the transformation feels predictable, responsive, and meaningful.

Body ownership refers to the experience that a body or body part is one's own. Ownership has often been studied through hand and body illusions, where visual, tactile, proprioceptive, and motor cues are combined. Synchronous and spatially plausible cues can support ownership, while incongruent cues can reduce it. In VR, ownership may be influenced by the visual appearance of the avatar, the spatial alignment between the physical and virtual body, and the correspondence between real and virtual movement.
% TODO cite: rubber hand illusion / virtual body ownership.
% TODO cite: avatar morphology / similarity literature, e.g. Tsakiris/Longo/Waltemate if used.
% TODO cite: multisensory congruence and body ownership.

In the context of rowing, body ownership is relevant because the user may see a virtual body performing the rowing action. An ergometer-congruent representation may place the virtual hands and body in closer correspondence with the felt body and handle. This could support ownership because the virtual body appears to occupy a plausible relation to the user's own body. A boat-centric representation may show a body engaged in a more recognizable rowing action, but if the visual hand or arm movement diverges from the felt handle movement, ownership may be affected differently.

Self-location concerns the experienced location of the self in relation to a body or environment. In immersive VR, the user's viewpoint is usually placed inside or near a virtual body, which can support the impression of being located in the virtual scene. Self-location may be influenced by visual perspective, body alignment, movement, and multisensory cues.
% TODO cite: self-location / full-body illusion literature.
% TODO cite: Lenggenhager et al. or related self-location studies.

Self-location is relevant to VR rowing because the user's physical body is stationary in a room, while the virtual scene may depict movement through water or a rowing environment. The user may feel located on the ergometer, in the virtual boat, or in a hybrid relation between both. In an ergometer-congruent mode, the felt machine and visible apparatus may reinforce the physical setup. In a boat-centric mode, the visual environment may encourage the user to interpret themselves as being in a boat, even though the physical cues still come from the ergometer.

The three components of embodiment therefore provide different lenses on the same mapping problem. Agency emphasizes the relation between intention, movement, and feedback. Ownership emphasizes whether the virtual body or body parts are experienced as belonging to the user. Self-location emphasizes where the user experiences themselves to be situated. A central assumption of this thesis is that the ergometer-congruent and boat-centric representations may produce different patterns across these components rather than a single uniform increase or decrease in embodiment.
